\chapter{IL TICCHETTIO DI OROLOGI COSTOSI}\label{il-ticchettio-di-orologi-costosi}

\hfill\break

L'elezione federale si avvicinava rapidamente. Jason intendeva usarla
come diversivo.

«Rimettimi in sesto» aveva detto. C'era un modo per farlo, insisteva.
Non era convenzionale. Non era approvato dalla FDA. Ma era una terapia
con una storia lunga e ben documentata. E mise in chiaro che intendeva
sfruttarla, provarci con o senza la mia collaborazione.

E visto che Molly lo aveva privato di quasi tutto ciò che era importante
per lui - e lasciato me tra le macerie - accettai di aiutarlo.
(Pensando, con ironia, a ciò che E.D. mi aveva detto anni prima: ``Mi
aspetto che ti prenda cura di lui. Mi aspetto che usi il buon senso.''
Era quello che stavo facendo?)

Nei giorni precedenti all'elezione di novembre Wun Ngo Wen ci informò
sulla procedura da seguire e sui rischi associati.

Consultarsi con Wun non era facile. Il problema non era tanto la rete di
sicurezza che lo circondava, malgrado fosse abbastanza difficile da
eludere, ma la massa di analisti e specialisti che si nutrivano dei suoi
archivi come i colibrì del nettare. Erano studiosi rinomati, controllati
dall'FBI e dalla Difesa Nazionale, che avevano giurato di mantenere il
segreto almeno pro tempore, ipnotizzati dalla vasta banca dati del
sapere marziano che Wun aveva portato con sé sulla Terra. Le
informazioni digitali ammontavano a più di cinquecento volumi, mille
pagine a volume, di astronomia, biologia, matematica, fisica, medicina,
storia e tecnologia, e la maggior parte erano di gran lunga più avanzate
delle conoscenze terrestri. Se fosse stato recuperato con la macchina
del tempo l'intero contenuto della Biblioteca di Alessandria,
difficilmente avrebbe potuto produrre un interesse più frenetico da
parte degli accademici.

Queste persone venivano messe sotto pressione per completare il loro
lavoro prima dell'annuncio ufficiale della presenza di Wun. Il governo
federale voleva una catalogazione almeno grossolana degli archivi (molti
erano scritti nella nostra lingua ma in modo approssimativo; alcuni,
invece, erano stati redatti utilizzando la scrittura scientifica
marziana) prima che anche i governi stranieri cominciassero a pretendere
di accedervi. Il Dipartimento di Stato programmava di produrre e
distribuire copie ripulite da cui fossero state espunte o ``presentate
in forma sintetica'' certe tecnologie potenzialmente preziose o
pericolose, mentre gli originali dovevano rimanere altamente riservati.

E così intere classi di studiosi lottavano e custodivano gelosamente il
loro contatto con Wun, che poteva interpretare o spiegare le lacune nel
testo marziano. In diverse occasioni fui scacciato dall'alloggio di Wun
da garbati uomini e donne del ``gruppo di fisica delle alte energie'' o
``del gruppo di biologia molecolare'' che, eccitati, pretendevano il
quarto d'ora accordato. Di tanto in tanto Wun mi presentava a queste
persone, ma nessuno era mai felice di vedermi, e quando Wun annunciò di
aver scelto me come suo medico personale la capogruppo delle scienze
mediche si agitò quasi al punto di farsi venire un infarto.

Jase rassicurò gli studiosi accennando al fatto che facevo parte del
``processo di socializzazione'' con cui Wun affinava i suoi modi
terrestri fuori dal contesto politico o scientifico, e io promisi alla
capogruppo che non avrei somministrato a Wun alcuna terapia senza il suo
coinvolgimento diretto. Tra i ricercatori si diffuse la voce che ero un
opportunista civile che si era fatto strada nella cerchia ristretta di
Wun manipolando gente a destra e sinistra, e che la mia ricompensa
sarebbe stata un corposo contratto per un libro pubblicabile una volta
che l'esistenza di Wun fosse stata divulgata. La voce nacque
spontaneamente ma non facemmo nulla per smentirla; serviva ai nostri
scopi.

L'accesso ai farmaci fu più facile di quanto mi aspettassi. Wun era
arrivato sulla Terra con un'intera farmacopea di medicine marziane di
cui non esisteva un corrispondente terrestre e un giorno, stando a
quanto affermava, avrebbe potuto averne bisogno per curare se stesso. Le
scorte di medicinali erano state confiscate dal veicolo con cui era
atterrato, ma gli erano state restituite nel momento in cui era stato
confermato il suo status di ambasciatore. (Senza dubbio il governo aveva
raccolto dei campioni, ma Wun dubitava che una semplice analisi avrebbe
rivelato lo scopo di quei materiali altamente sofisticati.) Fornì alcune
fiale del farmaco puro a Jason, che le portò fuori della Perielio
avvolto nella nube misteriosa dei privilegi da dirigente.

Wun mi diede istruzioni su dosaggio, tempi di somministrazione,
controindicazioni e potenziali problemi. Rimasi sgomento dalla lunga
lista dei rischi connessi. Anche su Marte, disse Wun, il tasso di
mortalità della transizione a Quarto era dello 0,1 per cento, un dato
non irrilevante, e il caso di Jason era complicato dalla SMA.

Ma senza la cura la prognosi di Jason era anche peggiore. E lui sarebbe
comunque andato avanti, con o senza la mia approvazione: in un certo
senso, il medico che prescriveva era Wun Ngo Wen, non io. Io avevo solo
il compito di supervisionare la procedura e trattare eventuali effetti
collaterali imprevisti. Il che mi alleggeriva la coscienza, anche se
sarebbe stato difficile difendere queste argomentazioni in
tribunale\ldots{} Wun poteva anche avere ``prescritto'' i farmaci, ma
non sarebbe stata sua la mano che li avrebbe iniettati nel corpo di
Jason.

Sarebbe stata la mia.

Wun Ngo Wen non sarebbe stato nemmeno presente. Jase aveva chiesto un
congedo di tre settimane tra la fine di novembre e i primi di dicembre,
e per allora Wun sarebbe diventato una celebrità globale, un nome (per
quanto insolito) riconosciuto da tutti. Wun sarebbe stato impegnato a
parlare alle Nazioni Unite e ad accettare l'ospitalità di un compendio
di monarchi, mullah, presidenti e primi ministri del nostro pianeta con
le mani un po' sporche di sangue, mentre Jason sudava e vomitava,
proiettato verso uno stato di salute migliore.

Avevamo bisogno di un posto dove andare. Un posto in cui potesse star
male senza dare nell'occhio, un posto in cui potessi assisterlo senza
attirare attenzioni indesiderate, ma abbastanza civilizzato da poter
chiamare un'ambulanza se le cose si fossero messe male. Un posto
confortevole. Un posto tranquillo.

«Conosco il luogo perfetto» disse Jason.

«Dov'è?»

«La Grande Casa.»

Mi misi a ridere finché non mi resi conto che diceva sul serio.

ooo

Diane non richiamò che una settimana dopo la visita di Lomax alla
Perielio, una settimana dopo che Molly aveva lasciato la città per
reclamare il compenso che le era stato promesso da E.D. Lawton o da uno
dei suoi scagnozzi.

Domenica pomeriggio. Ero solo nel mio appartamento. Una giornata
soleggiata, ma le tende erano tirate. Dopo una settimana passata a
dividermi tra i pazienti della clinica della Perielio e le lezioni
segrete con Wun e Jase, contemplavo la desolazione di quel vuoto fine
settimana. Era bello essere impegnati, ragionavo tra me e me, perché
quando sei impegnato sei immerso nei problemi quotidiani, innumerevoli
ma comprensibili, che tolgono spazio al dolore e soffocano il rimorso.
Era salutare. Era un modo di reagire. O perlomeno una tattica dilatoria.
Utile ma, ahimè, temporanea. Perché prima o poi il rumore si
affievolisce, la folla si disperde e torni a casa tua, con la lampadina
bruciata, la stanza vuota e il letto sfatto.

Stavo abbastanza male. Non sapevo neanche esattamente come mi sentivo -
o meglio, quale dei tanti dolori contrastanti e incompatibili avrei
dovuto accettare per primo. «Stai meglio senza di lei» mi aveva detto
Jase un paio di volte, ed era tanto vero quanto scontato: meglio senza
di lei, ma ancora meglio se fossi riuscito a trovare un senso, se fossi
riuscito a capire se Molly mi aveva usato o punito per averla usata, se
il mio amore gelido e forse un po' falso fosse stato equivalente al
ripudio freddo e proficuo di quell'amore, da parte sua.

Poi squillò il telefono, e fu imbarazzante perché ero impegnato a
togliere le lenzuola dal letto, appallottolandole per portarle nella
lavanderia, tanti detersivi e acqua bollente per cancellare l'aura di
Molly. Non vuoi essere interrotto durante una simile attività. Ti fa
sentire un pochino a disagio. Ma ero sempre stato schiavo di un telefono
che squilla. Risposi.

«Tyler?» disse Diane. «Sei tu, Ty? Sei solo?»

Ammisi di essere solo.

«Bene, sono contenta di averti trovato, finalmente. Volevo dirti che
stiamo cambiando numero di telefono. Non lo faccio mettere nell'elenco.
Ma nel caso avessi bisogno di metterti in contatto con me\ldots»

Mi dettò il numero privato che scrissi su un tovagliolo a portata di
mano. «Perché non lo fai mettere nell'elenco?» Lei e Simon avevano solo
il telefono fisso, ma immaginai che si trattasse di una penitenza
religiosa, come indossare abiti di lana o mangiare solo cereali
integrali.

«Da un lato perché riceviamo delle strane telefonate da E.D. Un paio di
volte ha chiamato di sera tardi e ha iniziato a fare la paternale a
Simon. Sembrava un po' ubriaco, francamente. E.D. odia Simon, E.D. lo ha
odiato dal primo giorno, ma dal trasferimento a Phoenix non l'avevamo
più sentito. Fino a ora. Il silenzio era doloroso, questo però è
peggio.»

Il numero di telefono di Diane poteva essere uno dei dati che Molly
aveva rubacchiato dal mio programma di gestione familiare e aveva
passato a E.D. Non potevo spiegarlo a Diane senza violare il giuramento
di sicurezza, per la stessa ragione per cui non potevo menzionare Wun
Ngo Wen o i replicatori mangiaghiaccio. Quello che le dissi fu che Jason
era stato coinvolto da suo padre in una lotta per il controllo della
Perielio uscendone vincitore, e che forse era questo che tormentava E.D.

«Potrebbe essere» ribatté Diane. «In effetti, così a ridosso del
divorzio\ldots»

«Quale divorzio? Parli di E.D. e Carol?»

«Jason non te l'ha detto? E.D. vive in affitto a Georgetown da maggio.
Le trattative sono ancora in corso, ma pare che a Carol andranno la
Grande Casa e gli alimenti e a E.D. tutto il resto. Il divorzio è stata
un'idea di lui, non di lei. È comprensibile, forse. Per decenni Carol è
stata a tanto così dal coma etilico. Non era un granché come madre e non
può essere stata un granché come moglie.»

«Stai dicendo che approvi?»

«Per niente. Non ho cambiato idea su di lui. Era un genitore terribile,
indifferente\ldots{} almeno per me. Non mi piaceva e non gli importava
piacermi. Ma non ne avevo neanche soggezione, non come Jason. Jason lo
vedeva come un enorme re dell'industria, un eccelso potente di
Washington\ldots»

«Non lo è?»

«Ha successo e una certa influenza, ma queste cose sono tutte relative,
Ty. Ci sono diecimila E.D. Lawton in questo paese. E.D. non sarebbe mai
arrivato da nessuna parte, se suo padre e suo zio non avessero
finanziato la sua prima attività\ldots{} sono certa che per loro non sia
stata altro che una spesa detraibile dalle tasse, niente di più. E.D.
era bravo in ciò che faceva, e quando lo Spin ha presentato
un'opportunità ne ha approfittato attirando l'attenzione di persone
veramente potenti. Ma, fondamentalmente, per i pezzi grossi ha
continuato a essere un arricchito. Non ha mai avuto dalla sua quelle
cose da Yale, Harvard e confraternite segrete e potenti. Per me niente
ballo delle debuttanti. Noi eravamo i bambini poveri del quartiere.
Voglio dire, era un bel quartiere, ma ci sono i ricchi di famiglia e ci
sono gli arricchiti, e noi sicuramente eravamo gli arricchiti.»

«Direi che le cose sembravano diverse, dall'altra parte del prato. Come
se la cava Carol?»

«La medicina di Carol sgorga dalla solita bottiglia di sempre. E tu,
invece? Come va tra te e Molly?»

«Molly è andata» dissi.

«Andata tipo ``andata al negozio'' o\ldots»

«Andata e basta. Ci siamo lasciati. Non trovo un eufemismo carino per
dirlo.»

«Mi spiace, Tyler.»

«Grazie, ma è la cosa migliore. Lo dicono tutti.»

«Tra me e Simon va tutto bene» disse sebbene non gliel'avessi chiesto.
«Questa faccenda della chiesa però è pesante per lui.»

«Questioni di ``politica'' di chiesa?»

«Il Jordan Tabernacle ha qualche problema giudiziario. Non conosco tutti
i dettagli. Non siamo coinvolti direttamente, ma Simon la sta prendendo
piuttosto male. Sicuro di stare bene, comunque? Ti sento un po' rauco.»

«Sopravvivrò.»

ooo

La mattina prima delle elezioni preparai un paio di valigie (vestiti
puliti, qualche libro tascabile, il mio kit medico) e passai a prendere
Jason a casa sua per andare in Virginia. Jase era ancora appassionato di
auto di lusso, ma dovevamo viaggiare senza attirare l'attenzione. Con la
mia Honda, quindi, non con la sua Porsche. Di quei tempi le autostrade
interstatali non erano sicure per viaggiare in Porsche.

La presidenza di Garland era stata propizia per chiunque avesse un
reddito superiore a mezzo milione di dollari, avversa per tutti gli
altri. Ed era piuttosto evidente dall'aspetto della strada, un paesaggio
scorrevole di magazzini al dettaglio incuneati fra centri commerciali
abbandonati e sbarrati, parcheggi in cui gli abusivi vivevano in auto
senza ruote, città autostradali che si sostentavano con i guadagni
provenienti da un autogrill e un autovelox. Cartelli di avvertimento
piazzati dalla polizia statale annunciavano {VIETATO FERMARSI DOPO IL
	TRAMONTO} oppure {CHIAMARE IL 911 DA UN NUMERO IDENTIFICABILE PER PRONTO
	INTERVENTO D'EMERGENZA}. La pirateria autostradale aveva dimezzato il
volume del traffico dei piccoli veicoli. Passammo gran parte del viaggio
incastrati fra autoarticolati a doppio rimorchio, alcuni visibilmente
mal ridotti, e automezzi verde mimetico per il trasporto delle truppe,
diretti a varie basi militari.

Ma non ne parlammo. E non parlammo neanche di elezioni, il cui risultato
era comunque già scontato, con Lomax che avrebbe ottenuto più voti di
ciascun candidato, dei due principali e dei tre rivali minori. Né di
replicatori mangiaghiaccio o di Wun Ngo Wen e certamente non di E.D.
Lawton. Parlammo, invece, dei vecchi tempi e di buoni libri e per gran
parte del tempo non parlammo affatto. Avevo riempito la memoria
dell'autoradio con quel genere di jazz spigoloso e controcorrente che a
Jason piaceva: Charlie Parker, Thelonious Monk, Sonny Rollins\ldots{}
gente che molto tempo prima aveva compreso la distanza fra la strada e
le stelle.

Ci fermammo davanti alla Grande Casa al tramonto.

La casa era illuminata a giorno, grandi finestre di un giallo color
burro sotto un cielo d'inchiostro iridescente. Quell'anno il clima in
periodo elettorale era rigido. Carol Lawton scese dal portico per
venirci incontro, il corpo minuto avvolto in una sciarpa di cachemire e
un maglione lavorato ai ferri. Era quasi sobria, a giudicare dalla sua
andatura ferma seppure un po' troppo artefatta.

Jason scese lentamente, con prudenza, dal sedile passeggero.

Era in fase di remissione o ci si avvicinava, per quanto in quei giorni
fosse possibile. Con un piccolo sforzo poteva sembrare normale. La cosa
che mi sorprese fu che smise di sforzarsi di apparire tale non appena
arrivammo alla Grande Casa. Attraversò l'ingresso sbandando, dirigendosi
verso la sala da pranzo. Non c'erano domestici - Carol aveva disposto
che avessimo la casa tutta per noi per un paio di settimane - ma il
cuoco aveva lasciato un vassoio di carne fredda e verdure nel caso
arrivassimo affamati. Jason si lasciò cadere su una sedia.

Io e Carol lo raggiungemmo. Era visibilmente invecchiata, dalla morte di
mia madre. I capelli erano così sottili che si intravedeva il contorno
del cranio, rosa e scimmiesco, e quando le presi il braccio ebbi la
sensazione di toccare un ramoscello secco rivestito di seta. Le guance
erano incavate. Gli occhi avevano l'iperattività fragile e nervosa di un
bevitore in astinenza. Quando le dissi che era bello rivederla, sorrise
mestamente. «Grazie, Tyler. So che ho un aspetto orribile. Gloria
Swanson in \emph{Viale del tramonto}. Non sono ancora pronta per un
primo piano, ti ringrazio di cuore, cacchio.» Non sapevo di cosa stesse
parlando. «Ma sopravvivo. Come sta Jason?»

«Come sempre» risposi.

«Sei dolce a tergiversare. Ma lo so\ldots{} be', non posso dire di
sapere tutto. Però so che è malato. Questo mi ha detto. E so che si
aspetta di essere curato da te. Una cura non convenzionale ma efficace.»
Scansò il braccio e mi guardò negli occhi. «È \emph{efficace}, vero,
questo farmaco che hai intenzione di dargli?»

Ero troppo sorpreso per dire qualcos'altro se non un «Sì.»

«Perché mi ha fatto promettere di non fare domande. Immagino che vada
tutto bene. Jason si fida di te. Perciò, mi fido di te. Anche se quando
ti guardo, non posso fare a meno di vedere il bambino che vive nella
casa dall'altra parte del prato. Ma vedo un bambino anche quando guardo
Jason. Bambini scomparsi\ldots{} non riesco a capire dove li ho persi.»

ooo

Quella notte dormii in una stanza per gli ospiti della Grande Casa, una
camera che avevo visto solo di sfuggita dal corridoio negli anni in cui
avevo vissuto nella proprietà.

Comunque, dormii solo in parte. Parte della notte la passai a letto
sveglio, cercando di valutare i rischi legali che mi ero assunto andando
là. Non sapevo esattamente quali leggi o protocolli avesse potuto
violare Jase sottraendo in modo illecito preparati farmacologici
marziani dal campus della Perielio, ma mi ero già reso complice del
fatto.

Il mattino dopo Jason si chiese dove avremmo dovuto conservare le varie
fiale di liquido trasparente che Wun gli aveva fornito, sufficienti a
curare quattro o cinque persone. («Nel caso perdessimo una valigia»
aveva spiegato all'inizio del viaggio. «Ridondanza.»)

«Ti aspetti una perquisizione?» Mi immaginai i funzionari federali con
le tute di protezione da rischio biologico salire le scale della Grande
Casa.

«Certo che no. Ma non è mai una cattiva idea tutelarsi dai rischi.» Mi
diede un'occhiata più da vicino, anche se gli occhi si muovevano a
scatti verso sinistra a intervalli di qualche secondo, un altro sintomo
della malattia. «Sei un po' in apprensione?»

Dissi che potevamo nascondere le scorte nella casa dall'altra parte del
prato, a meno che non dovessero stare al freddo.

«Secondo Wun sono chimicamente stabili in qualsiasi condizione a meno di
una guerra termonucleare. Ma un mandato di perquisizione includerebbe
tutta la proprietà.»

«Non so nulla di mandati, però so dove sono i nascondigli.»

«Mostrameli» disse Jason.

E così attraversammo il prato in fila, Jason dietro di me un po'
instabile. Era pomeriggio presto, il giorno delle elezioni, ma nello
spazio erboso tra le due case avrebbe potuto essere un qualsiasi
autunno, un qualsiasi anno. Da qualche parte nella macchia boschiva
attraversata dal torrente un uccello annunciò la sua presenza, un'unica
nota che cominciò audace ma andò affievolendosi come se avesse cambiato
idea. Poi raggiungemmo la casa di mia madre, girai la chiave e aprii la
porta su un'immobilità più profonda.

La casa era stata pulita e spolverata periodicamente, ma in sostanza era
chiusa dalla morte di mia madre. Non ero mai tornato a sistemare le sue
cose, non c'erano altri parenti e Carol aveva preferito lasciare
l'abitazione com'era piuttosto che cambiarla. Ma non era immortale.
Tutt'altro. Il tempo vi aveva nidificato. Il tempo si era messo comodo.
Il salotto puzzava di chiuso, di odori emanati dalla tappezzeria
indisturbata, carta gialla, tessuto deciso. In inverno, come Carol mi
disse in seguito, la casa veniva mantenuta calda a sufficienza da
evitare che i tubi congelassero; in estate le tende venivano chiuse per
ripararla dal caldo. Quel giorno faceva freddo, dentro e fuori.

Jason attraversò la soglia tremando. La sua andatura era stata
traballante tutta la mattina, per questo mi aveva permesso di portare i
farmaci (a parte quelli che avevo già messo da parte per la sua cura),
mezzo chilo o quasi di vetro e biochimici in una ventiquattrore di pelle
imbottita di gommapiuma.

«È la prima volta che vengo qui,» disse timidamente, «da prima che
morisse. È stupido dire che mi manca?»

«No, non è stupido.»

«È stata la prima persona che mi abbia trattato in maniera gentile.
Tutta l'amabilità della Grande Casa entrava dalla porta con Belinda
Dupree.»

Lo accompagnai, passando per la cucina, alla piccola porta che conduceva
al seminterrato. La casetta sulla proprietà dei Lawton era stata
progettata per assomigliare a un cottage del New England, o all'idea che
se ne aveva; scendemmo nella cantina di cemento grezzo dal soffitto così
basso che Jason dovette chinarsi per seguirmi. Lo spazio era abbastanza
grande da contenere caldaia, scaldabagno, lavatrice e asciugatrice.
L'aria lì sotto era anche più fredda e aveva un odore umido, minerale.

Mi accovacciai nell'angolino dietro la struttura in lamiera della
caldaia, uno di quegli anfratti polverosi che anche gli addetti alle
pulizie professionali ignorano abitualmente. Spiegai a Jase che lì c'era
una soletta di cartongesso rotta e con un po' di destrezza la si poteva
sollevare e scoprire l'esiguo spazio, non isolato, fra le travi di pino
e le fondamenta.

«Interessante» disse Jason dalla sua posizione a un metro dietro di me,
di fianco alla caldaia inattiva. «Che tenevi là dentro, Tyler? Vecchi
numeri di \emph{Gent}?»

Quando avevo dieci anni avevo conservato lì qualche giocattolo, non
perché temevo che qualcuno li rubasse ma perché era divertente sapere
che erano nascosti e che solo io potevo trovarli. In seguito avrei
nascosto cose meno innocenti: tentativi brevi di tenere un diario,
lettere per Diane mai consegnate o neanche finite e, sì, anche se non lo
avrei confessato a Jason, pornografia relativamente innocua stampata da
Internet. Tutti quei segreti colpevoli erano stati eliminati da tempo.

«Avremmo dovuto portare una torcia» brontolò Jase. L'unica lampadina in
alto gettava una luce irrilevante in quell'angolo coperto di ragnatele.

«Ce n'era una sul tavolo vicino al portafusibili.» C'era ancora. Mi
ritrassi dall'intercapedine quanto bastava per prenderla dalla mano di
Jason. Emetteva il bagliore annacquato e pallido di una batteria
morente, ma funzionava abbastanza da riuscire a trovare il pezzo
allentato di cartongesso senza andare a tastoni. Lo sollevai e vi feci
scivolare dietro la valigetta, poi sistemai il pannello al suo posto e
sparsi la polvere di gesso sulle fessure visibili.

Ma prima di riuscire a indietreggiare e uscire feci cadere la torcia che
rotolò ancora più in fondo dietro la caldaia, nel buio infestato dai
ragni. Feci una smorfia e mi allungai per prenderla, seguendo lo
sfarfallio del bagliore. Ne toccai il manico. Toccai anche
qualcos'altro. Qualcosa di vuoto ma consistente. Una scatola.

La tirai avvicinandola.

«Hai finito là dentro, Ty?»

«Un attimo» risposi.

Puntai la luce sulla scatola. Era una scatola da scarpe. Una scatola da
scarpe con un polveroso logo New Balance e una scritta spessa in
inchiostro nero: {RICORDI} ({SCUOLA}).

Era la scatola che mancava di sopra dall'étagère di mia madre, quella
che non ero riuscito a trovare dopo il funerale.

«Problemi?» chiese Jason.

«No» risposi.

Potevo indagare in seguito. Spinsi la scatola dove l'avevo trovata e
strisciai fuori dallo spazio polveroso. Mi alzai e mi spazzolai le mani.
«Direi che qui abbiamo finito.»

«Ricordatelo per me» disse Jason. «Nel caso me ne dimenticassi.»

ooo

Quella sera guardammo i risultati delle elezioni sull'enorme ma datato
impianto video dei Lawton. Carol non trovava le lenti correttive e si
sedette vicino allo schermo, strizzando gli occhi. Aveva passato gran
parte della sua vita da adulta ignorando la politica - «È sempre stata
di competenza di E.D.» - e dovemmo spiegarle chi erano alcuni dei
candidati più importanti. Ma sembrava godersi la situazione. Jason
faceva battute carine e Carol lo compiaceva ridendo, e quando rideva
vedevo qualcosa di Diane sul suo viso.

Tuttavia, Carol si stancava facilmente, perciò nel momento in cui le
emittenti cominciarono a dare i risultati stato per stato era già andata
in camera sua. Non ci furono sorprese. Alla fine Lomax prese tutto il
Nord-Est e gran parte del Midwest e dell'Ovest. Andò meno bene nel Sud,
ma anche lì il voto contrario fu diviso quasi in parti uguali tra i
democratici della vecchia scuola e i conservatori cristiani.

Cominciammo a raccogliere le tazze di caffè nel momento in cui l'ultimo
candidato del partito avversario riconosceva la sconfitta con un
discorso mogio ma gentile.

«Quindi vincono i buoni» dissi.

Jase sorrise. «Non credo che ce ne fossero.»

«Pensavo che Lomax fosse l'uomo giusto per noi.»

«Forse. Ma non fare l'errore di pensare che Lomax si preoccupi della
Perielio o del programma dei replicatori, se non come modo conveniente
di ridurre le spese destinate allo spazio e facendolo apparire come un
grande salto in avanti. Il denaro federale che avanza sarà destinato al
bilancio militare. Ecco perché E.D. non è riuscito a trovare tra i suoi
vecchi compari dell'industria aerospaziale qualcuno che volesse opporsi
sul serio a Lomax. Lomax non lascerà morire di fame la Boeing o la
Lockheed Martin. Vuole solo che si riorganizzino.»

«Per la difesa» aggiunsi. La tregua nel conflitto mondiale seguita al
caos derivato dallo Spin era passata da tempo. Forse un
riequipaggiamento dell'esercito non era un'idea così sbagliata.

«Se credi a quel che dice Lomax.»

«Tu no?»

«Temo di non potermelo permettere.»

E con quella chiosa me ne andai a letto.

La mattina dopo gli somministrai la prima iniezione. Jason si sdraiò sul
divano nel grande salotto dei Lawton, di fronte alla finestra. Indossava
jeans e una camicia di cotone, aveva un'aria disinvolta da patrizio
romano, fragile ma a suo agio. Se era spaventato non lo dava a vedere.
Si arrotolò la manica destra per scoprire l'incavo del gomito.

Presi una siringa dal kit, ci attaccai un ago sterile e la riempii con
il liquido chiaro di una delle fiale che non avevamo occultato.

Wun mi aveva assistito durante le prove. I protocolli della Quarta Età.
Se fossimo stati su Marte, ci sarebbe stata una cerimonia tranquilla e
un ambiente rassicurante. Lì ci dovemmo accontentare della luce di
novembre e del ticchettio di orologi costosi.

Gli strofinai la pelle prima dell'iniezione. «Non devi guardare.»

«Ma voglio farlo. Fammi vedere com'è fatto.»

Gli era sempre piaciuto sapere come funzionavano le cose.

ooo

L'iniezione non produsse effetti immediati, ma verso mezzogiorno del
giorno successivo aveva qualche linea di febbre.

Per quanto lo riguardava, disse, era come un leggero raffreddore e a
metà pomeriggio mi pregò di prendere termometro e manicotto della
pressione e\ldots{} be', di portarli da un'altra parte. Il concetto era
questo.

Allora mi sollevai il colletto per ripararmi dalla pioggia (una pioggia
vuota, scioccamente ostinata iniziata la notte e che persisteva nel
pomeriggio) e ancora una volta attraversai il prato verso casa di mia
madre, dove recuperai la scatola {RICORDI} ({SCUOLA}) dal seminterrato e
la portai di sopra nel salotto.

Dalle tende entrava la luce offuscata dalla pioggia. Accesi una lampada.

Mia madre era morta all'età di cinquantasei anni. Per diciotto anni
avevo condiviso quella casa con lei. Poco più di un terzo della sua
vita. Dei restanti due terzi avevo visto solo ciò che aveva scelto di
mostrarmi. Di tanto in tanto mi aveva parlato di Bingham, la sua città
natale. Sapevo, per esempio, che aveva vissuto con il padre (un agente
immobiliare) e la matrigna (un'assistente all'infanzia) in una casa in
cima a un ripido viale alberato; che aveva avuto un'amica d'infanzia di
nome Monica Lee; che c'era stato un ponte coperto, un fiume chiamato il
Little Wyecliffe e una chiesa presbiteriana che aveva smesso di
frequentare dopo aver compiuto sedici anni e in cui non era più entrata
se non per i funerali dei suoi genitori. Ma non aveva mai menzionato
l'università di Berkeley o ciò che aveva sperato di realizzare con la
sua laurea in economia aziendale o la ragione per cui aveva sposato mio
padre.

Una o due volte aveva preso le scatole per mostrarmene il contenuto, per
farmi comprendere che era sopravvissuta agli anni impossibili precedenti
alla mia nascita. Quelle erano le sue prove, Reperti A, B e C, tre
scatole di {RICORDI} e {CIANFRUSAGLIE}. Ripiegati in quelle scatole, da
qualche parte, c'erano frammenti di storia vera, verificabile: le prime
pagine color mou dei giornali che annunciavano attacchi terroristici,
guerre dichiarate, presidenti eletti o incriminati. Anche qui c'erano i
ninnoli che mi piaceva tenere in mano da bambino. Un pezzo da cinquanta
centesimi ossidato coniato nell'anno di nascita di suo padre; quattro
conchiglie rosa e marrone chiaro della spiaggia di Cobscook Bay.

{RICORDI} ({SCUOLA}) era la scatola che mi piaceva di meno. Conteneva
una spilla elettorale di qualche candidato democratico a un'alta carica
che evidentemente non ce l'aveva fatta, che mi piaceva per i suoi colori
brillanti, ma il resto dello spazio era occupato dal suo diploma,
qualche pagina strappata dal suo annuario delle superiori e un fascio di
piccole buste che non avevo mai voluto (o non mi era stato permesso)
toccare.

In quel momento ne aprii una e ne lessi il contenuto, quel tanto che
bastò per capire di cosa si trattasse: a) era una lettera d'amore e b)
la scrittura era ben diversa da quella ordinata di mio padre, nelle
missive in {RICORDI} ({MARCUS}).

Quindi mia madre aveva avuto un fidanzato al college. Era una notizia
che avrebbe potuto scombussolare Marcus Dupree (lo aveva sposato una
settimana dopo la laurea) ma difficilmente avrebbe scioccato qualcun
altro. Di sicuro non c'era motivo di nascondere la scatola nel
seminterrato, visto che era stata in bella vista per anni e anni.

Ma poi, era stata davvero mia madre a nasconderla? Non sapevo chi
potesse essere entrato in casa tra il suo ictus e il mio arrivo il
giorno dopo. Era stata Carol a trovarla esanime sul divano e forse in
seguito qualcuno del personale della Grande Casa aveva aiutato a
ripulire, e là dentro doveva esserci stato qualche paramedico che
l'aveva preparata per il trasporto. Ma nessuno di questi avrebbe avuto
un motivo neanche lontanamente plausibile per portare di sotto {RICORDI}
({SCUOLA}) e farlo scivolare nello spazio buio tra la caldaia e il muro
del seminterrato.

Comunque non era importante. Qualcuno l'aveva spostata, ma non era stato
commesso nessun crimine. Poteva benissimo esser stato il poltergeist
della casa. Con ogni probabilità non l'avrei mai saputo e non aveva
senso soffermarsi sulla questione. Tutto di quella stanza, ogni oggetto
della casa, incluse quelle scatole, prima o poi avrebbe dovuto essere
recuperato, venduto o buttato via. L'avevo sempre rimandato, Carol
l'aveva sempre rimandato, ma quel lavoro doveva esser già stato fatto.

Intanto, però\ldots{}

Intanto, rimisi {RICORDI} ({SCUOLA}) sull'ultima mensola dell'etagere
tra {RICORDI} (M{ARCUS}) e {CIANFRUSAGLIE}. E svuotai completamente la
stanza.

ooo

L'obiezione medica più preoccupante che avevo sollevato con Wun Ngo Wen
riguardo alla cura di Jason era la questione delle interazioni con altri
farmaci. Non potevo interrompere la sua terapia convenzionale senza
causargli una disastrosa ricaduta. Mi turbava altrettanto, però,
combinare il suo regime farmacologico quotidiano con la ristrutturazione
biochimica di Wun.

Wun mi aveva assicurato che non ci sarebbero stati problemi. Il
trattamento della longevità non era un ``farmaco'' nel senso
convenzionale. Ciò che stavo iniettando nel sangue di Jason era più
simile a un programma informatico attivato biologicamente. I farmaci
convenzionali in genere interagiscono con proteine e superfici
cellulari. La pozione di Wun interagiva direttamente con il DNA.

Tuttavia doveva entrare in una cellula per svolgere il suo compito e
superare nel percorso la chimica del sangue di Jason e il suo sistema
immunitario per arrivare a destinazione, giusto? Wun aveva risposto con
enfasi che questo non aveva importanza. Il cocktail della longevità era
abbastanza flessibile da agire in qualsiasi condizione fisiologica
tranne che in quella di morte.

Il gene che causava la SMA, però, non era mai migrato sul pianeta rosso
e là i farmaci che prendeva Jase erano sconosciuti. E malgrado Wun
avesse insistito che le mie preoccupazioni erano infondate, notavo che
nel dirlo sorrideva di rado. Perciò cercammo di mitigare i rischi. Avevo
iniziato a ridurre le medicine per la SMA di Jason nella settimana
precedente alla prima iniezione. Ma non le avevo sospese del tutto, solo
diminuite.

Era sembrato che la strategia funzionasse. Nel momento in cui eravamo
arrivati alla Grande Casa, con quel carico di farmaci più leggero Jason
mostrava solo sintomi minori, perciò cominciammo il trattamento con
ottimismo.

Tre giorni dopo aveva accessi di febbre alta che non riuscivo a far
scendere. Il giorno dopo era già quasi sempre in stato di semicoscienza.
Il giorno dopo ancora la pelle gli era diventata rossa e aveva iniziato
a coprirsi di vesciche. Quella sera cominciò a urlare.

Continuava a urlare nonostante la morfina che gli somministravo.

Non era un urlo a piena gola, ma un gemito che periodicamente saliva di
volume, un suono che ti aspetteresti da un cane malato, non da un essere
umano. Era del tutto involontario. Quando era lucido non emetteva alcun
suono né si ricordava di averlo fatto, sebbene si ritrovasse la laringe
infiammata e dolorante.

Carol diede una gran prova di coraggio nel sopportare la situazione. In
alcune zone della casa il lamento di Jason era quasi impercettibile - le
stanze da letto sul retro, la cucina - e lei passava lì gran parte del
tempo, leggendo o ascoltando la radio locale. Ma la tensione era
evidente e in breve tempo ricominciò a bere.

Forse non dovrei dire ``ricominciò''. Non aveva mai smesso. In realtà lo
aveva ridotto a quel minimo che le permettesse di connettere,
barcamenandosi tra il terrore puro dell'improvvisa astinenza e la
tentazione di una completa ebbrezza. E spero di non risultare
superficiale: Carol stava percorrendo un sentiero difficile, ed era
riuscita a rimanervi così a lungo per amore di suo figlio, per quanto
quell'amore in tutti quegli anni fosse rimasto latente. A farla
deragliare era stato il suono del dolore di Jason.

All'inizio della seconda settimana dall'avvio del procedimento cominciai
a somministrare a Jase delle endovenose tenendo d'occhio la pressione in
aumento. Era stato abbastanza bene per l'intera giornata, nonostante
l'aspetto terribile, coperto di croste dove la pelle viva non era
esposta, gli occhi quasi sepolti nella carne gonfia circostante. E
lucido a sufficienza da chiedere se Wun Ngo Wen avesse già fatto la sua
prima apparizione in televisione. (Non ancora. Era prevista per la
settimana successiva.) Ma verso l'imbrunire era ricaduto in uno stato di
incoscienza e il lamento, assente da un paio di giorni, ricominciò, a
piena gola e doloroso da sentire.

Doloroso per Carol, che comparve alla porta della stanza da letto con le
guance rigate dalle lacrime e un'espressione di rabbia feroce e vetrosa.
«Tyler» disse «devi farlo smettere!»

«Faccio quel che posso. Non risponde agli oppiacei. Sarebbe meglio
parlarne domattina.»

«Non lo \emph{senti}?»

«Certo che lo sento.»

«Non significa niente? Quel suono non significa \emph{niente} per te?
Cazzo! Sarebbe stato meglio con un ciarlatano del Messico. Sarebbe stato
meglio con un santone. Hai la \emph{minima} idea di cosa gli hai
iniettato? Stramaledetto ciarlatano!»

Purtroppo ripeteva domande che avevo già preso a pormi io stesso. No,
\emph{non} sapevo cosa gli stessi iniettando, non in senso rigoroso e
scientifico. Avevo creduto alle promesse dell'uomo di Marte, ma non
sarebbe stata una giustificazione sufficiente per Carol. Il procedimento
in sé era più difficile, di certo più angosciante di quanto avessi
previsto. Forse stava funzionando male. Forse non stava funzionando
affatto.

Jase emise un urlo lugubre che terminò in un sospiro. Carol si mise le
mani sulle orecchie. «Sta \emph{soffrendo}, ciarlatano del cazzo!
\emph{Guardalo}!»

«Carol\ldots»

«Niente Carol, macellaio! Chiamo un'ambulanza. Chiamo la polizia!»

Le andai incontro e la presi per le spalle. La sentivo fragile ma
pericolosamente viva sotto le mani, un animale in trappola. «Carol,
ascoltami.»

«Perché, perché dovrei \emph{ascoltarti}!»

«Perché tuo figlio ha messo la sua vita nelle mie mani. Ascolta. Carol,
ascolta. Avrò bisogno di qualcuno che mi aiuti. Vado avanti da giorni
senza dormire. Fra non molto avrò bisogno di qualcuno che stia con lui,
qualcuno che sia davvero ferrato in medicina e che possa prendere
decisioni fondate.»

«Avresti dovuto portare un'infermiera.»

Avrei dovuto, ma non mi era stato possibile, e comunque non era quello
il punto. «Non ho un'infermiera. Dovresti farlo tu.»

Le ci volle un po' per recepirlo. «Io!»

«A quanto ne so, sei ancora abilitata per esercitare la professione.»

«Non esercito da\ldots{} forse decenni? Decenni\ldots»

«Non ti sto chiedendo di operare al cuore. Voglio solo che tu tenga
d'occhio pressione e temperatura. Puoi farlo?»

La rabbia si dissipò. Ne era lusingata, spaventata. Ci penso su. Poi mi
rivolse uno sguardo duro. «Perché dovrei aiutarti? Perché dovrei essere
complice di questa \emph{tortura}?»

Stavo ancora formulando una risposta da darle quando una voce alle mie
spalle disse: «Oh, vi prego.»

Era la voce di Jason. Una delle caratteristiche della terapia con il
farmaco marziano era lo stato di lucidità che andava e veniva a caso, a
suo piacimento. Evidentemente gli era appena tornata. Mi voltai.

Fece una smorfia e tentò, con poco successo, di mettersi seduto. Ma
aveva gli occhi vigili.

Si rivolse a sua madre: «Sul serio, non ti pare un po' fuori luogo? Per
favore, fa' quello che ti chiede Tyler. Sa quello che fa, e anch'io.»

Carol lo fissava. «Ma io non\ldots{} non ho\ldots{} cioè, non
posso\ldots»

Poi si voltò e uscì dalla stanza in modo malfermo, una mano poggiata al
muro.

Mi sedetti accanto a Jase. La mattina dopo Carol entrò nella stanza,
sembrava umiliata ma sobria e si offrì di sostituirmi. Jason era
tranquillo e in realtà non aveva bisogno di essere accudito, ma lasciai
che se ne occupasse lei e andai a recuperare un po' di sonno.

Dormii per dodici ore. Quando tornai nella camera, Carol era ancora lì,
teneva la mano del figlio privo di sensi e gli carezzava la fronte con
una tenerezza che non le avevo mai visto prima.

ooo

La fase di ripresa cominciò dopo una settimana e mezzo dall'inizio del
trattamento. Non ci fu nessuna transizione improvvisa, nessun momento
magico. I momenti di lucidità cominciarono a protrarsi e la pressione si
stabilizzò quasi nell'intervallo di normalità.

La sera del discorso di Wun alle Nazioni Unite trovai un televisore
portatile nella zona della casa riservata ai domestici e lo portai nella
camera di Jason. Carol ci raggiunse poco prima che iniziasse la
trasmissione.

Non credo che Carol credesse a Wun Ngo Wen.

La sua presenza sulla Terra era stata annunciata ufficialmente il
mercoledì precedente. Ormai la sua immagine spiccava sulle prime pagine
da giorni, oltre alle riprese in diretta di lui che percorreva il prato
della Casa Bianca sotto il braccio benevolo del presidente in carica. La
Casa Bianca aveva precisato che Wun si trovava lì per aiutarli ma che
non aveva una soluzione immediata al problema dello Spin né informazioni
nuove sugli Ipotetici. La reazione pubblica era stata cauta.

Quella sera, Wun montò sul palco della camera del Consiglio di Sicurezza
e salì sul podio che era stato adattato alla sua altezza. «Ma dai, è
solo un esserino minuscolo» disse Carol.

Jason intervenne: «Abbi un po' di rispetto. Rappresenta un'unica cultura
ininterrotta, durata più di una qualsiasi delle nostre.»

«Mi sembra che rappresenti più la corporazione dei lecca-lecca» rispose
Carol.

La sua dignità venne ripristinata nei primi piani. Alla telecamera
piacevano i suoi occhi e il suo sorriso sfuggente. E quando parlò al
microfono lo fece in modo sommesso, il che abbassò il tono effettivo
della sua voce a un livello più terrestre.

Wun sapeva (o era stato preparato a capire) quanto l'evento potesse
sembrare inverosimile al terrestre medio. (Nella sua presentazione, il
segretario generale aveva detto: «Viviamo davvero nell'epoca dei
miracoli.») Poi, con il suo miglior accento della regione atlantica
centrale, ci ringraziò tutti per la nostra ospitalità e raccontò con
nostalgia della sua casa e spiegò perché l'aveva lasciata per venire
sulla Terra. Dipinse Marte come un luogo sconosciuto ma del tutto umano,
quel genere di posto che sarebbe bello visitare, dove la gente era
amichevole e il paesaggio interessante, anche se gli inverni, ammise,
spesso erano rigidi.

(«Sembra il Canada» disse Carol.)

Quindi arrivò al nocciolo della questione. Tutti volevano sapere degli
Ipotetici. Purtroppo, il popolo di Wun ne sapeva su di loro, poco più di
noi\ldots{} gli Ipotetici avevano incapsulato Marte durante il suo
viaggio verso la Terra e i marziani erano inermi come lo eravamo stati
noialtri.

Non era in grado di immaginare quali fossero le motivazioni degli
Ipotetici. Da secoli si dibatteva sulla questione, ma anche i più grandi
pensatori marziani non l'avevano mai risolta. Era interessante, disse
Wun, che la Terra e Marte fossero stati isolati quando si trovavano
sull'orlo di catastrofi globali: «La nostra popolazione, come la vostra,
si avvicina al limite di sostenibilità. Sulla Terra l'industria e
l'agricoltura consumano petrolio, risorsa che va rapidamente
esaurendosi. Su Marte non abbiamo petrolio, ma dipendiamo da un'altra
materia prima che scarseggia, l'azoto elementare: sostiene il nostro
ciclo agricolo e impone limiti assoluti al numero di vite umane che il
pianeta può sostenere. Ce la siamo cavata un po' meglio di voi, ma solo
perché siamo stati costretti a riconoscere il problema sin dall'inizio
della nostra civiltà. Entrambi i nostri pianeti si trovavano e si
trovano di fronte alla possibilità di un collasso economico e agricolo e
a una mortalità umana catastrofica. Entrambi i pianeti sono stati
incapsulati prima di arrivare a quel punto.

Forse gli Ipotetici hanno compreso questa nostra realtà e magari questo
ha influenzato le loro azioni. Ma non lo sappiamo con assoluta certezza.
Né sappiamo cosa si aspettano da noi, semmai dovessero aspettarsi
qualcosa, né quando o se lo Spin finirà. Non \emph{possiamo} saperlo
finché non raccogliamo informazioni più dirette sugli Ipotetici.»

«Fortunatamente,» proseguì Wun mentre la camera lo inquadrava più da
vicino, «c'è un modo per raccogliere quelle informazioni. Sono venuto
qui con una proposta che ho discusso sia con il Presidente Garland che
con il neoeletto Lomax, e anche con altri capi di Stato» e seguitò
delineando le basi del progetto dei replicatori. «Se siamo fortunati, ci
dirà se gli Ipotetici hanno colpito altri mondi, come questi mondi hanno
reagito e quale potrebbe essere il destino finale della Terra.»

Ma quando cominciò a parlare della Nube di Oort e della tecnologia di
retroazione autocatalitica vidi gli occhi di Carol diventare vitrei.

«Non può essere vero» mormorò dopo che Wun aveva lasciato il podio di
fronte a gente che applaudiva sbalordita e gli opinionisti della rete
avevano cominciato a masticare e rigurgitare il suo discorso. Sembrava
davvero spaventata. «C'è qualcosa di vero in tutto questo, Jason?»

«Quasi tutto» confermò tranquillo Jason. «Non posso confermare quanto
siano rigidi gli inverni su Marte.»

«Siamo davvero sull'orlo del disastro?»

«Siamo sull'orlo del disastro da quando sono sparite le stelle.»

«Intendo per il petrolio e il resto. Se non ci fosse stato lo Spin,
staremmo tutti morendo di fame?»

«La gente muore \emph{già} di fame. Muore di fame perché non possiamo
mantenere sette miliardi di persone secondo la prosperità del modello
nordamericano senza spremere il pianeta come un limone. È difficile
sostenere quei numeri. Sì, è vero. Se lo Spin non ci uccide prima,
presto o tardi assisteremo a un avvizzimento globale dell'umanità.»

«E questo c'entra qualcosa con lo Spin?»

«Forse, ma né io né il marziano in televisione lo sappiamo per certo.»

«Ti stai prendendo gioco di me.»

«No.»

«Sì, invece. Ma non fa niente. Lo so che sono ignorante. Sono anni che
non sfoglio un giornale. C'era sempre il rischio di vedere la faccia di
tuo padre, tanto per dirne una. E l'unica cosa che guardo in TV sono i
telefilm del pomeriggio. In quelli non ci sono marziani. Mi sembra di
essere Rip van Winkle. Ho dormito troppo. E non mi piace molto il mondo
in cui mi sono risvegliata. Le uniche cose che non sono spaventose
sono\ldots» gesticolò verso la TV, «sono ridicole.»

«Siamo tutti Rip van Winkle» disse Jason dolcemente. «Tutti stiamo
aspettando di svegliarci.»

ooo

L'umore di Carol migliorava di pari passo con la salute di Jason, perciò
iniziò a mostrare un interesse più vivo per la prognosi di suo figlio.
La aggiornai sulla SMA, una malattia che ai tempi della sua laurea in
medicina non era stata ancora diagnosticata formalmente; fu un modo per
eludere le domande sulla cura in sé, un tacito accordo che sembrava
comprendere e accettare. La cosa importante era che la pelle devastata
di Jason stava guarendo e i campioni di sangue che avevo mandato a un
laboratorio di analisi di D.C. mostravano una riduzione drastica delle
placche proteiche neuronali.

Comunque, era ancora riluttante a parlare dello Spin e sembrava
scontenta quando io e Jase ne discutevamo in sua presenza. Ripensai alla
poesia di Housman che Diane mi aveva insegnato tantissimi anni prima:
``\emph{All'infante non è occorso / che se l'è mangiato l'orso}.''

Carol era stata attaccata da diversi orsi, qualcuno grosso quanto lo
Spin e qualcun altro piccolo quanto una molecola di etanolo. Credo che
avrebbe invidiato l'infante.

ooo

Diane chiamò (sul mio telefono personale, non su quello di casa di
Carol) qualche sera dopo l'apparizione di Wun alle N.U. Mi ero ritirato
nella mia stanza e Carol vegliava su Jason. Per tutto novembre la
pioggia era andata e venuta; in quel momento pioveva, la finestra della
camera era uno specchio fluido di luce gialla.

«Sei alla Grande Casa» disse Diane.

«Hai parlato con Carol?»

«La chiamo una volta al mese. Sono una figlia rispettosa. A volte è
abbastanza sobria da poterci parlare. Che ha Jason?»

«È una lunga storia. Sta migliorando. Non c'è nulla di cui
preoccuparsi.»

«Odio quando la gente dice così.»

«Lo so, ma è vero. C'era un problema, ma l'abbiamo risolto.»

«E puoi dirmi solo questo.»

«Per ora sì. Come vanno le cose a te e Simon?» L'ultima volta che
avevamo parlato mi aveva accennato a problemi legali.

«Non tanto bene. Ci trasferiamo.»

«Vi trasferite dove?»

«Fuori Phoenix, comunque. Lontano dalla città. Il Jordan Tabernacle per
ora è stato chiuso\ldots{} pensavo che forse ne avessi sentito parlare.»

«No» dissi - perché avrei dovuto sentir parlare dei problemi finanziari
di una chiesetta della Tribolazione del Sud-Ovest? - e continuammo a
discutere di altre questioni; Diane promise di darmi notizie quando lei
e Simon avessero avuto il nuovo indirizzo. Certo, perché no, che
diamine!

Ne sentii parlare, però, la sera dopo.

Stranamente, Carol insistette nel voler vedere l'ultimo notiziario.
Jason era stanco ma vigile e ben disposto, perciò tutti e tre resistemmo
fino in fondo a guardare quaranta minuti di minacce di intervento armato
internazionale e processi di celebrità. Qualche notizia era
interessante: c'era un aggiornamento su Wun Ngo Wen, che si trovava a un
incontro in Belgio con i funzionari dell'U.E. e buone notizie
dall'Uzbekistan, dove la base avanzata navale era stata finalmente
smantellata. Poi c'era un servizio speciale sulla SDC e l'industria
lattiera israeliana.

Guardammo immagini drammatiche di bestiame abbattuto che veniva gettato
in fosse comuni e cosparso di calce. Cinque anni prima, l'industria
giapponese delle carni bovine aveva subito una devastazione simile. Era
scoppiata un'epidemia di bovini e ungulati colpiti da SDC, soppressa in
decine di paesi dal Brasile all'Etiopia. L'equivalente umano
dell'infezione era curabile con i moderni antibiotici ma restava un
problema scottante per le economie del terzo mondo.

I produttori israeliani di latte applicavano rigidi protocolli sulla
sepsi e le analisi, perciò nessuno si aspettava un'epidemia proprio lì.
Ma la cosa peggiore era che il caso primario - la prima infezione - era
stato rintracciato in una spedizione non autorizzata di ovuli fecondati
proveniente dagli Stati Uniti.

La spedizione fu ricondotta a un'organizzazione benefica tribolazionista
chiamata ``Parola per il Mondo'', la cui sede si trovava in una zona
industriale fuori Cincinnati, Ohio. Perché la PpiM contrabbandava ovuli
di bestiame in Israele? Venne fuori che non era per motivi
particolarmente caritatevoli. Gli investigatori cominciarono le indagini
seguendo i patrocinatori della PpiM attraverso una decina di società
anonime fino a un consorzio di chiese tribolazioniste e
dispensazionaliste e gruppi politici marginali piccoli e grandi. Un
punto condiviso della dottrina biblica era l'interpretazione, comune in
questi gruppi, di un passo tratto dal Libro dei Numeri (capitolo
diciannove) e dedotta da altri versetti in Matteo e Timoteo, vale a dire
che la nascita in Israele di una giovenca rossa pura avrebbe segnato il
secondo Avvento di Gesù Cristo e l'inizio del Suo Regno sulla Terra.

Era un vecchio concetto. Gli estremisti ebrei affini credevano al
sacrificio di un vitello rosso sul Monte del Tempio che avrebbe
contrassegnato l'Avvento del Messia. Negli anni precedenti c'erano stati
diversi assalti da ``vitello rosso'' sulla Cupola della Roccia, uno dei
quali aveva danneggiato la moschea al-Aqsā e per poco non aveva
innescato una guerra regionale. Il governo israeliano aveva fatto del
suo meglio per reprimere il movimento, ma era riuscito solo a portarlo
alla clandestinità.

Secondo i notiziari, c'erano diversi caseifici patrocinati dalla PpiM in
tutte le regioni americane medio occidentali e del Sudovest, tutti
tranquillamente dediti al compito di accelerare l'Apocalisse. Avevano
cercato di ottenere una varietà di vitello puro rosso sangue,
presumibilmente superiore alle numerose giovenche deludenti presentate
come candidate nei quarant'anni precedenti.

Queste fattorie avevano eluso sistematicamente le ispezioni federali e i
protocolli di alimentazione animale, al punto da occultare un'epidemia
di SDC bovina che aveva attraversato il confine da Nogales. Gli ovuli
infetti generavano animali da riproduzione con moltissimi geni per il
mantello rosso, ma quando nascevano i vitelli (in un caseificio nel
Negev legato alla PpiM) morivano quasi tutti in età precoce per
insufficienza respiratoria. I cadaveri erano stati sepolti nel completo
silenzio, ma era già troppo tardi. L'infezione si era diffusa al
bestiame adulto e a un certo numero di allevatori.

Era una fonte di imbarazzo per l'amministrazione statunitense. La FDA
aveva già annunciato una revisione della sua politica e la Difesa
nazionale stava congelando i conti bancari della PpiM e consegnando
mandati a coloro che raccoglievano fondi per i tribolazionisti. Al
notiziario si vedevano immagini di agenti federali che uscivano da
edifici anonimi reggendo scatole di documenti e che mettevano lucchetti
alle porte di chiese oscure.

L'annunciatore citò il nome di alcune.

Una di quelle era il Jordan Tabernacle.